\section{Datasets \& Benchmarks}
\label{sec:app_bench}
% =============================================================================

This section describes the datasets and evaluation protocol for LLAMIA-Bench.
We find four major themes in how AI systems collaborate with domain-expert agents:
\emph{Behavioral imitation}: reproducing human play at a target skill level, the problem behind bots like Play Magnus\footnote{\url{https://www.playmagnus.com}} and Maia~\citep{mcilroy2020aligning}.
\emph{State assessment}: predicting human-aligned properties of game states such as difficulty and engagement, the core task in Lichess's puzzle rating system and Chess.com's adaptive training~\citep{lichess2024puzzles}.
\emph{Comparative}: explaining why a position favors one side---identifying material or structural advantages and disadvantages---as required in single-position analysis~\citep{jhamtani-etal-2018-learning} and game-level commentary.
\emph{Rationale}: generating natural-language explanations for why a player made a specific move, from single-move annotation~\citep{jhamtani-etal-2018-learning,zang2019automated} to the game-length narratives produced by channels like Agadmator and GothamChess.
These four themes place progressively harder demands on the communication channel between LLM and subagent: from a single-position state query (behavioral imitation) to game-length narrative integration (commentary), and from signals with partial textual correlates like move quality to ones without, such as aesthetic interest~(\cref{sec:app_bench_puzzle_int}).
LLAMIA-Bench instantiates each as an evaluation task.
Chess serves as the testbed because it offers all four at once: subagents whose internal representations are mapped by interpretability work~\citep{jenner2024evidence}, public game databases at scale~\citep{Lichess,Caissabase}, established benchmarks with dedicated task finetune baselines, and decades of human--engine collaboration.

% =============================================================================
\subsection{Dataset \& Metrics}
\label{sec:app_bench_dataset}
\label{sec:app_reward}
% =============================================================================

Table~\ref{tab:bench_dataset} summarizes dataset provenance;
Table~\ref{tab:bench_metrics} lists evaluation metrics per task.
Splits marked OOD are \emph{out-of-distribution}: the test distribution is absent or shifted relative to Stage-2 training, so generalization must come from internalized representations rather than memorization.
Where the original authors provide a fixed test split we use it; otherwise we sample a random held-out split.
Detailed descriptions of each task follow in \S\ref{sec:app_bench_tasks}.

%% ── TABLE 1: DATASET PROVENANCE ──────────────────────────────────────────────
\begin{table*}[ht]
\centering
\small
\setlength{\tabcolsep}{5pt}
\renewcommand{\arraystretch}{1.4}
\caption{\textbf{LLAMIA-Bench: dataset provenance.} \textbf{Ours} = constructed for this work. OOD = out-of-distribution test split.}
\label{tab:bench_dataset}
\vspace{4pt}
\begin{tabular}{p{4.8cm} p{5.8cm} p{2.4cm}}
\toprule
\textbf{Task / Source} & \textbf{Description} & \textbf{Train\,/\,Test} \\
\midrule
\makecell[tl]{Behavior Cloning (\S\ref{sec:app_bench_bc}) \\ {\scriptsize MAIA-KDD, Lichess~\cite{mcilroy2020aligning}}}
  & Predict the move a human of a given Elo would play; 5 Elo buckets (1100--1900)
  & 12M / authors' \\
\makecell[tl]{\quad\emph{In the wild} (OOD) \\ {\scriptsize MAIA-KDD, Lichess (\textbf{Ours})}}
  & Three OOD splits: GM-25 (top-25 GMs), Low-Time (clock pressure $<$10\%), Elo Gap ($>$500 pts)
  & $\leq$167K / --- \\
\midrule
\makecell[tl]{Puzzle Understanding (\S\ref{sec:app_bench_puzzle}) \\ {\scriptsize Lichess Puzzles~\cite{lichess2024puzzles}}}
  & Predict difficulty (Glicko-2) and interest (community votes) of tactical puzzles
  & 4M / 5K \\
\midrule
\makecell[tl]{Move Annotation (\S\ref{sec:app_bench_annotation}) \\ {\scriptsize Lichess~\cite{jhamtani-etal-2018-learning}}}
  & Generate natural-language explanation for a single move across 5 semantic categories
  & 90K / authors' \\
\midrule
\makecell[tl]{Game Commentary (\S\ref{sec:app_bench_commentary}) \\ {\scriptsize Agadmator YouTube (\textbf{Ours})}}
  & Produce coherent multi-turn narrative spanning an entire game
  & 1.9K / 100 \\
\bottomrule
\end{tabular}
\end{table*}

%% ── TABLE 2: METRICS ─────────────────────────────────────────────────────────
\begin{table*}[ht]
\centering
\small
\setlength{\tabcolsep}{5pt}
\renewcommand{\arraystretch}{1.3}
\caption{LLAMIA-Bench: evaluation metrics and RL reward signals. $\uparrow$ higher is better; $\downarrow$ lower is better. Section references point to detailed metric definitions.}
\label{tab:bench_metrics}
\vspace{4pt}
\begin{tabular}{p{3.5cm} p{5.2cm} p{3.4cm}}
\toprule
\textbf{Task} & \textbf{Metric(s)} & \textbf{RL Reward} \\
\midrule
Behavior Cloning       & Move-match accuracy $\uparrow$                                              & Top-3 rank \\
Puzzle Understanding   & Spearman $\rho$ $\uparrow$ (\S\ref{sec:app_bench_puzzle})                   & Normalized MAE $\downarrow$ \\
{\scriptsize\quad\emph{Difficulty}} & {\scriptsize Spearman $\rho$ $\uparrow$ (\S\ref{sec:app_bench_puzzle_diff})} & {\scriptsize Normalized MAE $\downarrow$} \\
{\scriptsize\quad\emph{Interest}}   & {\scriptsize Spearman $\rho$ $\uparrow$ (\S\ref{sec:app_bench_puzzle_int})}  & {\scriptsize Normalized MAE $\downarrow$} \\
{\scriptsize\quad\emph{Solved (\%)}} & {\scriptsize Exact solution-line accuracy (\S\ref{sec:app_bench_puzzle_solve}); parity metric, not primary} & {\scriptsize ---} \\
Move Annotation        & G-eval $\uparrow$;\enspace BLEU-2 $\uparrow$ (\S\ref{sec:app_bench_annotation}) & G-eval \\
Game Commentary        & G-eval $\uparrow$;\enspace BLEU-2 $\uparrow$ (\S\ref{sec:app_bench_commentary}) & G-eval \\
\bottomrule
\end{tabular}
\end{table*}

% =============================================================================
\subsection{Detailed Task Descriptions}
\label{sec:app_bench_tasks}
% =============================================================================

% ── GAME-LEVEL COMMENTARY ────────────────────────────────────────────────────
\subsubsection{Game-Level Commentary}
\label{sec:app_bench_commentary}

Game-level commentary requires a coherent, multi-turn narrative spanning an entire game---unlike move-level annotation, errors compound across the narrative, and the model must track evolving themes (initiative shifts, pawn-structure transformations, time trouble). We introduce \textbf{Agadmator-2K}, the first large-scale dataset for this task: 1,900 narrated games from Agadmator's YouTube channel,\footnote{\url{https://www.youtube.com/@agadmator}} totaling approximately 500 hours.

\paragraph{Dataset construction.}
Move-segmented commentary is unavailable from YouTube. We construct it in four steps: (i)~transcripts are extracted via Whisper-v3-large; (ii)~video timestamps are aligned to PGN move sequences using a sliding-window move-tracking buffer; (iii)~GPT-4o labels which moves each transcript segment references, guided by the known PGN; (iv)~segments are accepted only when the inferred move order matches the PGN exactly, discarding retries and out-of-order narration. The test set consists of 100 games, held out by ascending view count to minimize overlap with LLM pretraining corpora. This is a heuristic proxy for low contamination, not a guarantee; we discuss contamination further in \S\ref{sec:app_bench_contamination}.

\paragraph{Metrics.}
We use the same G-eval framework as move annotation (\S\ref{sec:app_bench_annotation}), adapted for game-level commentary: each generated segment is scored on relevance, completeness, clarity, and fluency, with the judge grounded by Lc0-BT4 engine lines and the ground-truth transcript. We also use the BLEU-2 metric. The RL reward is G-eval.

% ── BEHAVIOR CLONING ─────────────────────────────────────────────────────────
\subsubsection{Behavior Cloning}
\label{sec:app_bench_bc}

Behavior cloning measures whether LLAMIA can imitate human play conditioned on skill level or player identity. The evaluation metric across all splits is \textbf{move-match accuracy}: the fraction of positions where the model's top-1 predicted move exactly matches the target player's move. We follow the MAIA evaluation protocol~\citep{mcilroy2020aligning}: Maia variants use best-of-$N$ sampling; LLAMIA uses a single forward pass.

The RL reward uses a softer signal: the \textbf{top-3 rank} of the target move in the model's output distribution, normalized to $[0,1]$. Top-1 exact match as a reward collapsed training---the signal was too sparse for most positions, yielding near-zero gradients throughout Stage~2. Rank within the top-3 provides a dense, monotone reward that penalizes misranking without requiring exact prediction, while remaining consistent with the evaluation objective.

\paragraph{Maia Benchmark (Elo Buckets).}
\label{sec:app_bench_bc_maia}
We evaluate on the MAIA-KDD held-out test set~\citep{mcilroy2020aligning}, stratified into five Elo buckets: 1100, 1300, 1500, 1700, and 1900. The test set is player--game disjoint from all training data. Each bucket is treated as an independent task; the aggregate BC score reported in the main paper is the unweighted average across buckets.

\paragraph{GM-25 (OOD).}
\label{sec:app_bench_bc_gm}
GM-25 targets the top-25 rated grandmasters in FIDE history by peak rating.\footnote{\url{https://en.wikipedia.org/wiki/List_of_chess_players_by_peak_FIDE_rating}} Each grandmaster is a separate behavioral target. The largest available per-GM corpus is 4,641 games (Viktor Korchnoi),\footnote{\url{https://www.365chess.com/top-chess-players-games.php}} less than 3\% of the data that per-GM Maia models require~\citep{mcilroy2020aligning}. No Stage-2 training data is drawn from these GM corpora; generalization must come from internalized representations and cross-Elo behavioral transfer.

\paragraph{Low-Time (OOD).}
\label{sec:app_bench_bc_time}
Under severe clock pressure, players shift strategy regardless of position quality. We extract positions where either player's remaining clock is below 10\% of the initial time control, or where cumulative time usage differs by more than 50\% between the two sides. Positions are stratified by game phase (opening, middlegame, endgame) and sampled equally across time controls, yielding 129,000 positions. Clock-context metadata is absent from Stage-2 training, making this an OOD split: the model must infer time-pressure effects from the position and move alone.

\paragraph{Elo Gap (OOD).}
\label{sec:app_bench_bc_elo}
Players adapt their style when facing a large skill gap---weaker players take more risks, stronger players simplify. We filter Lichess Rapid and Classical games where the Elo difference exceeds 500 points, yielding 34,000 games (68,000 player-side instances). Extreme skill-gap matchups are rare in the Stage-2 training distribution; conditioning on opponent strength must emerge from contextual signals rather than memorization.

% ── PUZZLE UNDERSTANDING ─────────────────────────────────────────────────────
\subsubsection{Puzzle Understanding}
\label{sec:app_bench_puzzle}

Puzzle understanding probes whether LLAMIA has internalized the subagent's positional representations well enough to predict human-aligned properties of game states. Both sub-tasks draw from the same 4-million-puzzle Lichess corpus,\footnote{\url{https://database.lichess.org/lichess_db_puzzle.csv.zst}} which provides community-derived ground-truth labels for difficulty and engagement. We hold out a shared test set of 5,000 puzzles, stratified by difficulty (Glicko-2 quintiles), theme (tactical motif), and interest (score quintiles) to ensure uniform coverage across the label space. Evaluation uses Spearman~$\rho$ between predicted and ground-truth values; the RL reward is a normalized mean-absolute-error penalty.

\paragraph{Difficulty Estimation.}
\label{sec:app_bench_puzzle_diff}
Puzzle difficulty is operationalized via a Glicko-2 rating system~\citep{glickman2012example}: each human solving attempt is treated as a rated match between solver and puzzle, and the Glicko-2 rating accumulated over all attempts serves as ground truth. The model receives the puzzle position and solution line, and predicts difficulty on a normalized scale. Evaluation uses Spearman~$\rho$ between predicted values and ground-truth Glicko-2 ratings on the stratified 5,000-puzzle test split.

\paragraph{Interest Estimation.}
\label{sec:app_bench_puzzle_int}
Lichess assigns each puzzle an interestingness score (range: $-100$ to $+100$) computed from community upvotes and downvotes, weighted by solver performance. This signal has no straightforward textual correlate: a puzzle's aesthetic appeal depends on motif rarity, surprise, and solution elegance---features encoded in the subagent's positional representation but absent from any verbalized move list. The model predicts interest from the same input as difficulty; evaluation uses Spearman~$\rho$ on the same stratified 5,000-puzzle test split. Interest estimation is the diagnostic task on LLAMIA-Bench: the non-verbalizable nature of the target signal means that all text-mediated systems collapse on this task (\cref{sec:results_vd}).

\paragraph{Puzzle Solving Accuracy.}
\label{sec:app_bench_puzzle_solve}
We also report \textbf{Solved (\%)}: the fraction of test puzzles for which the model produces the complete correct solution line---every forced move in sequence---using policy-only decoding (single forward pass per position, no search). The model receives the initial puzzle FEN and outputs moves one at a time; a puzzle is marked solved only if all moves in the ground-truth solution are produced in the correct order.

This metric is excluded for GPT-5.1\,+\,Lc0 (marked~--- in \cref{tab:puzzle_understanding_extended}) because verbalized engine access makes it uninterpretable: a system that queries Lc0 at each puzzle position and forwards the top-ranked move would score near-perfect not by reasoning about the position but by delegating each step to the engine. The metric is informative only when the model must solve the puzzle from its own internalized representations without live tool queries. All other systems in Table~\ref{tab:puzzle_understanding_extended} use policy-only decoding for this column.

Puzzle solving accuracy functions as a \emph{parity metric} on LLAMIA-Bench: all systems with Lc0 access cluster in the 84--94\% range, and LLAMIA's improvement over LLAMIA-Verb is modest (2--3 pp). The metric confirms that engine-access systems are not deficient tactically; the differentiation between LLAMIA and LLAMIA-Verb arises in difficulty and interest prediction, not in puzzle-solving throughput.

% ── MOVE ANNOTATION ──────────────────────────────────────────────────────────
\subsubsection{Move Annotation}
\label{sec:app_bench_annotation}

Move annotation evaluates LLAMIA's ability to generate natural-language explanations for individual moves, conditioned on the board state and the move played. We follow the benchmark of \citet{jhamtani-etal-2018-learning}: 90,000 Lichess games annotated in English, with annotations categorized into five semantic dimensions that span both explanation themes from \S\ref{sec:app_bench}. The \emph{rationale} theme is instantiated by three dimensions---\emph{description} (what the move does), \emph{quality} (blunder, inaccuracy, good, best), and \emph{planning} (lookahead and intent)---while the \emph{comparative} theme is instantiated by two---\emph{context} (positional advantages and disadvantages relative to prior or future moves) and \emph{comparative} (alternative moves and why they were rejected). The standard benchmark provides the target move as input; we additionally evaluate zero-shot without this prior to test whether internalized representations can identify annotation-worthy moves.

\paragraph{Metrics.}
\textbf{BLEU-2} and perplexity~\citep{lee2022improving} are evaluated per annotation category, following prior work. The primary metric is \textbf{G-eval}~\citep{liu2023g}: an LLM-as-judge framework in which GPT-4o rates each generated annotation on a 0--1 scale across four dimensions (relevance, accuracy, completeness, fluency). The judge receives the board FEN, the move in algebraic notation, and Lc0-BT4's top-3 engine lines as grounding context, so its assessments are anchored in engine analysis rather than surface plausibility alone. Per-annotation G-eval scores are averaged across the four dimensions; the corpus-level score is the mean over all test annotations. G-eval also serves as the RL reward signal for this task.

% =============================================================================
\subsection{Data Contamination Statement}
\label{sec:app_bench_contamination}
% =============================================================================

All evaluation in LLAMIA-Bench is conditioned on board positions represented as FEN strings. We enforce a strict \textbf{FEN-level disjointness} guarantee: no FEN appearing in any test split co-occurs in any stage of training---projector pretraining (Stage~1), RL training (Stage~2), or the base LLM's supervised fine-tuning data. Concretely, we collect the set of all FENs used across projector pretraining pairs and Stage-2 RL rollouts, and verify that the intersection with each test split is empty. For the Maia BC test set, this property is inherited from the player--game disjoint split of \citet{mcilroy2020aligning}. For the puzzle understanding test split, the 5,000 held-out puzzles are sampled after removing all FENs present in the training pool. For Agadmator-2K, the 100 held-out games are additionally sorted by ascending view count as a heuristic to reduce overlap with LLM pretraining corpora, though we cannot verify disjointness with respect to closed-source pretraining data. For the OOD behavior-cloning splits (GM-25, Low-Time, Elo Gap), no Stage-2 training data is drawn from these distributions by construction; we further verify that no test FEN appears in the Stage-1 projector data.